\section*{الفصل \ref{ch:b2:palaeoenvironments} --- الحفريات والبيئات القديمة}

\begin{solution}{exo:b2:palaeoenvironments:1}
(أ) $112/135 = 0.830$ و$112/150 = 0.747$. (ب) تتجمّع الأصداف الميتة على القاع على مدى سنوات
أو أكثر: فهي تضمّ أنواعًا عاشت هناك من قبل، أو نادرًا، أو نُقلت إليه، فيكون
\omterm{def:b2:palaeoenvironments:assemblages}{تجمّع الموت} متوسّطًا لأجيال كثيرة من المجتمع. (ج) المجموعات رخوة الجسم، ولا
سيّما آكلات الفتات واللواحم كالديدان: فهي بلا هيكل تتحلّل قبل الدفن.
\end{solution}

\begin{solution}{exo:b2:palaeoenvironments:2}
(أ) $5 \times 165 = \qty{825}{kyr}$، أي نحو \qty{0.8}{Myr}، ضمن المدّة من 0.7 إلى مليون سنة.
(ب) $183.73 - 182.77 = \qty{0.96}{Myr}$، و$0.96/0.165 = 5.8$: نحو ستّة آفاق. (ج)
تتغيّر السلالة في أقلّ بكثير من مليون سنة؛ ولا تفصل بين الأشكال المتعاقبة إلا الآفاق،
في حين يجمعها الحيّز معًا.
\end{solution}

\begin{solution}{exo:b2:palaeoenvironments:3}
(أ) سحنة Nereites الأثرية: قاع البحر العميق في نظام عكرات، وراء حافة الرصيف، حيث
الغذاء شحيح. (ب) سحنة Skolithos الأثرية: الشاطئ الأمامي أو الوجه الشاطئي. (ج) سحنة Cruziana الأثرية:
بعيدًا عن الشاطئ، بين قاعدة أمواج الطقس الهادئ وقاعدة أمواج العواصف. (د) سحنة Zoophycos الأثرية: في
الميزوزوي، الرصيف السفلي أو المنحدر، تحت قاعدة أمواج العواصف.
\end{solution}

\begin{solution}{exo:b2:palaeoenvironments:4}
(أ) تحوي الطبقة A مجموعات ضيّقة تحمّل الملوحة: ماء بحر عادي. والطبقة B، الفقيرة بالأنواع والغنية
بالأفراد، مع فتات النباتات ودون أيّ من المجموعات ضيّقة تحمّل الملوحة، توحي بماء قليل الملوحة،
كمصبّ نهري مثلًا. (ب) المجموعات الغائبة عديدة، شائعة وجيّدة الحفظ في الطبقات البحرية
من العمر نفسه: فغيابها المشترك نمط، في حين قد تغيب حفرية واحدة مصادفةً أو
لسوء الحفظ. (ج) الراسب (قنوات مدّية أو نهرية، وأغشية طينية، وجذور)، والحفريات
الأثرية (سحنة Cruziana أثرية قليلة الأنواع)، وانتقال الطبقة B، جانبيًّا، إلى
رواسب بحرية أو نهرية.
\end{solution}

\begin{solution}{exo:b2:palaeoenvironments:5}
(أ) $P = 1/2$ حين $kh/s = \ln 2$: $s = 0.05 \times 0.1/0.693 = \qty{0.0072}{m.yr^{-1}}$،
أي \qty{7.2}{mm} في السنة. (ب) تنصيف $s$ يضاعف الأسّ: $P = (1/2)^2 = 0.25$، أي النصف.
(ج) تتنصّف $k$: $P = e^{-2.5} = 0.082$ بدلًا من $e^{-5} = 0.0067$، أي أرجح بمقدار $e^{2.5} = 12$
مرّة. (د) الأصداف السميكة المتينة ممثّلة بأكثر من حصّتها، وكلّما كان الدفن أبطأ كان ذلك أشدّ.
\end{solution}

\begin{solution}{exo:b2:palaeoenvironments:6}
(أ) الحدّ الأعلى $183.73 + 0.35 = \qty{184.08}{Ma}$: فقيمة الجدول \qty{184.2}{Ma} تقع
خارجه، بفارق \qty{0.12}{Myr}. (ب) $183.73 - 182.77 = \qty{0.96}{Myr}$. (ج) $6/0.96 =
\qty{6.25}{m}$ لكل مليون سنة. (د) يُعرَّف الحيّز الأحيائي بالحفريات، التي يكاد ظهورها
واختفاؤها يتزامنان على مساحات واسعة؛ أمّا سماكة الصخر المترسّب في
ذلك الزمن فتتوقّف على معدّل التراكم المحلّي.
\end{solution}

\begin{solution}{exo:b2:palaeoenvironments:7}
(أ) $0.49\,e^{-1046/3704} = 0.369$ و$0.49\,e^{-1166/3704} = 0.358$؛ والمتوسط $0.364$. (ب) $120
\times (1 - 0.364) = \qty{76.3}{m}$. (ج) $76.3/(1 - 0.49) = \qty{149.6}{m}$. (د) يفترض العامل
أن للطبقة المستعادة كلّها مسامية السطح $\varphi_0$؛ لكنّ طبقة
سماكتها \qty{150}{m} عند السطح أقلّ مسامية عند قاعدتها منها عند قمّتها، فتملأ الحبيبات نفسها
أقلّ من \qty{149.6}{m}.
\end{solution}

\begin{solution}{exo:b2:palaeoenvironments:8}
(أ) يحوي حجم الوحدة $\varphi$ من الماء و$1 - \varphi$ من الحبيبات: $\rho = \varphi
\rho_{\text{ماء}} + (1 - \varphi)\rho_{\text{حبيبة}}$، ومنه الصيغة؛ و$(2710 -
2573)/(2710 - 1000) = 0.080$. (ب) $0.51\,e^{-956/4545} = 0.413$. (ج) التلاحم: فالكالسيت
المترسّب في المسام، الذي تخضع له الكربونات إلى جانب الانضغاط، يزيل من المسامية أكثر بكثير
ممّا يزيله الدفن وحده.
\end{solution}

\begin{solution}{exo:b2:palaeoenvironments:9}
(أ) عكرات البحر العميق (سحنة Nereites الأثرية)؛ وقاع جوراسي للرصيف السفلي أو المنحدر
(Zoophycos)؛ وقاع بعيد عن الشاطئ بين قاعدتي الأمواج (Cruziana)؛ ووجه شاطئي (Skolithos)؛
ومستنقع على اليابسة. (ب) ضحل الماء نحو الأعلى: تحرّك الساحل نحو البحر. (ج) انخفاض مستوى البحر،
أو إمداد بالرواسب أسرع من الهبوط، كما أمام دلتا نامية؛ والانخفاض
يوجد في الزمن نفسه في أحواض أخرى، أمّا الدلتا فمحلّية ورواسبها تأتي
من جهة واحدة. (د) في تتابع متوافق تتعاقب البيئات المتجاورة؛ والقفز
فوق البيئات الوسيطة يدلّ على فجوة في الزمن أو على صدع.
\end{solution}

\begin{solution}{exo:b2:palaeoenvironments:10}
(أ) زمن الدفن $h/s$ = 100 و10 و\qty{1}{yr}. الجسم الرخو: $e^{-100} \approx
\num{4e-44}$، و$e^{-10} = \num{4.5e-5}$، و$e^{-1} = 0.37$. والصدفة ($k = \qty{0.01}{yr^{-1}}$):
$e^{-1} = 0.37$، و$e^{-0.1} = 0.90$، و$e^{-0.01} = 0.99$. (ب) $s = kh/\ln 100 = 0.1/4.605 =
\qty{0.022}{m.yr^{-1}}$، أي \qty{22}{mm} في السنة. (ج) بالكاد، على مدى أشهر؛ لكنّ طبقة عاصفة أو
عكرة أو سقوط رماد يدفن القاع في ساعات. (د) ثبات معدّل الترسيب: فيكون
زمن الدفن عندئذ أقصر بكثير من $h/s$ بالمعدّل المتوسط، و$P = e^{-kT}$، حيث $T$
زمن الدفن الفعلي، يؤول إلى الواحد حين يكون $T$ أصغر بكثير من $1/k$.
\end{solution}

\begin{solution}{exo:b2:palaeoenvironments:11}
(أ) تحوي طبقة رقيقة $\mathrm dz$ على العمق $z$ الحبيبات $(1 - \varphi(z))\,\mathrm dz$؛
وهي عند السطح تملأ $\mathrm dz_0$ مع $(1 - \varphi_0)\,\mathrm dz_0 = (1 -
\varphi(z))\,\mathrm dz$. (ب) $\varphi = 0.63\,e^{-816/1960} = 0.415$؛ و$(1 - 0.415)/0.37 =
1.581$، أي \qty{158.1}{m}. (ج) أكثر ممّا ينبغي بمقدار \qty{9.5}{m}: فكما في \cref{exo:b2:palaeoenvironments:7}،
ليست الطبقة المستعادة كلّها بمسامية سطحها. (د) معدّل الحبيبات،
$h(1-\varphi)/\Delta t$، لا يستعمل إلا السماكة والمسامية الحاليتين، وكلتاهما مقيسة.
\end{solution}

\begin{solution}{exo:b2:palaeoenvironments:12}
(أ) $9.5/0.165 = 57.6$ أفقًا. (ب) $180/57.6 = \qty{3.1}{m}$ من الحجر الجيري اليوم، و$338/57.6 =
\qty{5.9}{m}$ من الراسب الطازج. (ج) \qty{1.5}{m} على الأكثر: $180/1.5 = 120$ فاصلًا، و121
عيّنة. (د) انقطاعات الترسيب، وطبقات تغيب عنها الأمونيتات، وآفاق مكثّفة في
بضعة سنتيمترات حيث وصلت الرواسب ببطء.
\end{solution}

\begin{solution}{pb:b2:palaeoenvironments:1}
\textbf{1.} $866 \sin \ang{12} = 866 \times 0.2079 = \qty{180.0}{m}$.
\textbf{2.} $481 \times 0.2079 = \qty{100.0}{m}$.
\textbf{3.} الحجر الرملي، بمبدأ التراكب: فالجناحان لا يميلان إلا بزاوية \ang{12}، فالطبقات ليست
مقلوبة والأدنى هو الأقدم.
\textbf{4.} $866 - 766 = \qty{100}{m}$ و$1046 - 866 = \qty{180}{m}$.
\textbf{5.} \qty{766}{m} على الأقلّ هناك؛ وقمّته الأصلية محتوتة، فسماكته الكاملة
مجهولة.
\textbf{6.} 0.70 عند السطح؛ و$0.70\,e^{-766/1408} = 0.406$ عند \qty{766}{m}.
\textbf{7.} إنها مجموعات ضيّقة تحمّل الملوحة: كان للماء ملوحة البحر المفتوح.
\textbf{8.} سحنة Cruziana أثرية: المنطقة البعيدة عن الشاطئ، بين قاعدة أمواج الطقس الهادئ وقاعدة أمواج
العواصف.
\textbf{9.} العواصف: بلغت أمواجها القاع، فكسرت الأصداف ورسّبتها في طبقات
ذات قواعد حاتّة.
\textbf{10.} كان ماء القاع فقيرًا بالأكسجين: قليل من الحفّارات، وبقيت الصفائح، وتكوّن البيريت.
\textbf{11.} سحنة Zoophycos الأثرية، وهي في الميزوزوي الرصيف السفلي أو المنحدر، تحت
قاعدة أمواج العواصف؛ مع أكسجين كافٍ أحيانًا للحفّارات.
\textbf{12.} لا: فالأمونيتات سبحت في الماء فوقه؛ وغرقت أصدافها أو انجرفت إلى القاع
بعد الموت فشكّلت تجمّع موت.
\textbf{13.} $184.2 - 174.7 = \qty{9.5}{Myr}$.
\textbf{14.} $180/9.5 = \qty{18.9}{m}$ لكل مليون سنة.
\textbf{15.} $(2710 - 2573)/(2710 - 1000) = 0.080$.
\textbf{16.} $180 \times (1 - 0.080)/(1 - 0.51) = \qty{338}{m}$.
\textbf{17.} $338/9.5 = \qty{35.6}{m}$ لكل مليون سنة.
\textbf{18.} $180 \times 0.92/9.5 = \qty{17.4}{m}$ لكل مليون سنة.
\textbf{19.} أعطى الطين الصفحي 11.3 و\qty{4.4}{m} لكل مليون سنة: فانخفض معدّل الحبيبات
نحو أربع مرّات. وتوقّف وصول كربونات الضحضاح والرصيف مع تعمّق
القاع، ولم يبلغه إلا مطر بطيء من الطين.
\textbf{20.} تقدّم: أنهار، ثم ضحضاح، فقاع بعيد عن الشاطئ، فقاع تحت قاعدة أمواج
العواصف، ثم البحر المفتوح للطباشير. فالعمود يسجّل تجاوزًا بحريًّا.
\textbf{21.} $57 - 41 = \qty{16}{m}$ في أثناء الحجر الجيري، و$86 - 57 = \qty{29}{m}$ في أثناء الطين الصفحي،
و\qty{45}{m} في المجموع.
\textbf{22.} $(60 - 5) + 338 + 148.6 = \qty{542}{m}$، في $184.2 - 161.5 = \qty{22.7}{Myr}$:
أي \qty{23.9}{m} لكل مليون سنة.
\textbf{23.} $45/542 = 0.083$: نحو \qty{8}{\percent}، وأقلّ إن كان مخطّط الدورات يبالغ في
السعة. وجاء الباقي من هبوط القاع.
\textbf{24.} تراكم الطين الجيري لحجر فاردن الجيري بنحو
\textbf{\qty[text-series-to-math=true]{36}{m}} لكل مليون سنة حين ترسّب.
\end{solution}
